V1 C1 fuel flow at ICAO 85 \% thrust, kg/s  0.2858  0.2689  +6.3  5 & Engine & The C1 is modelled as a lower rating of the A2's hardware; the two variants' own databank rows differ by more than the limit at part power. \\
V5 TSFC at 45,000 ft, M 0.564, lb/h/lbf  0.5653  0.6225  -9.2  8 & Engine & A bare engine against the Citation X's installed one (bleed and power extraction); pyCycle's bare engine is 10.6 \% under the same figure. The installation factor on fuel comes from this difference. \\
V6b thrust at 43,835 ft at the fitted limit, lbf  1382  1335  +3.5  2 & Engine & After the fan-speed limit was fitted to the 45,000 ft row; the thesis's table is not exact to the 2 \% asked of it. \\
X0b A2 fuel flow at ICAO 85 \% thrust, lb/h  2876  3079  -6.6  5 & Engine against pyCycle & pyCycle's generic component maps at part power; the cycle model of the report is within 1.2 \% of the same row. \\
X1b C1 takeoff as a rating: pressure ratio  15.66  16.58  -5.6  5 & Engine against pyCycle & Generic compressor maps run at the C1's lower rating. \\
X2a TSFC at 1,500 ft, M 0.234, lb/h/lbf  0.4483  0.5309  -15.6  8 & Engine against pyCycle & A bare engine against an installed one, and generic maps at low altitude part power: about 10 \% of model uncertainty there. \\
X2b TSFC at 45,000 ft, M 0.564, lb/h/lbf  0.5562  0.6225  -10.6  8 & Engine against pyCycle & A bare engine against an installed one, as V5: both bare models are about 10 \% under the installed engine. \\
X3 fan corrected speed, 45,000 ft over 1,500 ft  1.462  1.258  +16.2  5 & Engine against pyCycle & pyCycle's generic fan map, run at 118 \% of its top speed line. \\
3-D model scaled by span: length, ft  measured  57.11  44.40  +28.6  1.0 & Geometry and mass & NASA's 3-D model is an artist's model, not to scale: scaled to the span it is 29 \% too long. It gives planform shape and widths only (its tail dihedral, 47.1 deg from the raw slope, was replaced by the published 50 deg that two head-on photographs confirm); lengths and heights come from NASA's side-view drawing. \\
3-D model scaled by span: height, ft  measured  20.78  15.20  +36.7  3.0 & Geometry and mass & As above: the 3-D model is not to scale. \\
3-D model scaled by span: wing area, ft2  measured  597.80  540.00  +10.7  3.0 & Geometry and mass & As above: the 3-D model is not to scale. \\
3-D model scaled by length: span, ft  measured  90.34  116.20  -22.3  1.0 & Geometry and mass & As above: scaled to the length its span is 22 \% short. \\
3-D model scaled by length: height, ft  measured  16.15  15.20  +6.2  3.0 & Geometry and mass & As above; the FlightGear model is squashed to the published 15.2 ft. \\
3-D model, span by span and chords by length: wing area, ft2  measured  464.70  540.00  -13.9  3.0 & Geometry and mass & As above; the wing's chords are a fit to the published 540 ft2. \\
side view: height, ft (NASA 15.2)  measured  14.34  15.20  -5.6  3.0 & Geometry and mass & NASA gives 15.2 ft and Northrop Grumman 14.6 ft; the drawing, scaled by the length, gives 14.3 ft, inside the limit of the second and outside that of the first. \\
wing as built: root chord with fillet against the side view, ft  measured  8.85  7.59  +16.5  5.0 & Geometry and mass & The wing as built keeps the 3-D model's root fillet (8.85 ft) where the drawing and photographs show 7.59 ft; with the drawing's chords the area falls short of the published 540 ft2. Open. \\
B3x fuel flow at it with 20 kW taken off, at 10000 ft, Mach 0.27: change, \%  -0.47178  -0.51432  +8.27\%  2.0 & JSBSim model as generated & The limit, 2 \% of the change, was written for changes of the size of the thrust's; the fuel flow's change at maximum thrust there nets two opposite effects out to -0.5 \%, and the 0.04 points of the fuel flow itself by which the tables' interpolation in Mach number differs from the cycle's are 8 \% of it. Kept as a miss; the limit was not changed. \\
B5 ref build trims at cruise & JSBSim model as generated & AeroSandbox's build-up does not carry a V-tail panel's twist into its pitching moment: 3 deg of tail incidence moves its moment by 0.0005, where the strip model's tail moves it by 0.11 (fdm/aero\_ref.py; a rotated section was tried and moved it by 0.008). The reference build therefore lacks the tail incidence the calibration set, and at the cruise needs 13 deg of ruddervator and nearly all the thrust. Its comparison with the strip model was made before and stands; the build is not used further. \\
S3 section maximum lift, Mach 0.6, Re 2 million  1.198  1.423  -15.8  10 & Section data and lift curve & The section data stops at 1.20, where its correction puts a shock stall; the tunnel reached 1.42. XFOIL cannot say: it does not converge past the first sonic point at Mach 0.6. A fit was made (fdm/cfd/stall\_fit.py: the tunnel's maximum lift with the CFD's shape) and met S3, but its blend back to the table's own deep stall after the CFD's last angle (10 to 12 deg) falls 0.26 to 0.37 per degree against the CFD's steepest 0.13, so by the rule written with it, it was withdrawn and the miss stays (docs/PLAN.md). The model stalls early at cruise Mach numbers, on the safe side for the laws' angle-of-attack limit. \\
A1 trimmed lift-curve slope, Mach 0.59 to 0.62, per deg  0.1341  0.1161  +15.5  10 & Section data and lift curve & From the added drag and not from a changed wing: the check trims six weights at 57,293 ft, and with the CFD's wave drag the three heaviest (21,254 lb and over) no longer hold height at full throttle, so the slope is fitted to the three lightest, 0.1341 against 0.1226 from four before. Taking the wave drag out brings it back to +5.6 \%; taking the stall fit or the tail angle out does not. It reflects a ceiling about 500 ft low at heavy weights (the mission's cruise-climb heights), inside their limit. \\
P1 time from 10,000 to 55,000 ft, s  4181.9  4595.6  -9.0 \%  5 \% & Climb, cruise-climb, endurance & The CFD's wave drag at Mach 0.55 to 0.7 was added to the section data and the drag increment refitted, 0.00586 to 0.00355 (below the 0.0040 to 0.0075 band written beforehand) and, after the tail's and the gear's lateral positions were measured, to 0.00308, fan-speed factor 1.01. One increment and one factor cannot hold the 40,000 to 55,000 ft climb rates (+11.1 and +15.6 \%), the cruise-climb heights and the total time together: the fit's compromise leaves the time 9.0 \% short against a limit of 5 (at 0.00355 the time passed, -0.1 \%, and the 737 minute height missed by 130 ft). No further fit: the limits allowed the one for each change (docs/PLAN.md, ''Wave drag from the CFD'', ''Lateral dimensions''). \\
A10 climb 10,000 to 55,000 ft hands off & Flight laws & The same climb time as the mission's P1: 4077 s flown by the laws against the quasi-steady 4452 s (-8.4 \%, limit 5), the refit's compromise between the 40,000 to 55,000 ft climb rates and the total time. It passed before the wave drag went in (4386 s, -1.5 \%). \\
O2 time from 10,000 to 55,000 ft (pressure) grows with the temperature & Off the design conditions: days, wind, turbulence & The limit was wrong, and is kept. Above FL220 the schedule is in Mach number (rising to 0.60 at 55,000 ft), so the climb rate goes with the true airspeed, which falls with the cold, while thrust and drag both go with pressure: the model's -20 K day is 3 \% slower (3,115 s against 3,021 s), and the hot days slower still (their service ceilings are about 55,000 ft at +15 K and 52,500 ft at +30 K). \\
W1 cruise in turbulence at 60,000 ft, severity 5 (sigma\_w 2.01 m/s): airspeed within 5 kt & Off the design conditions: days, wind, turbulence & The airspeed varies by the gust itself: the gust along the heading varies by 25 kt and the airspeed by 11, about half; no loop holds the airspeed within 5 kt in air that moves like this. The limit asked for what the air does not allow; height and bank, which the laws can hold, pass. \\
W1 cruise in turbulence at 60,000 ft, severity 5 (sigma\_w 2.01 m/s): load factor 0.7 to 1.3 & Off the design conditions: days, wind, turbulence & The lowest load factor in the severity-5 turbulence at 60,000 ft is 0.70 at the limit's edge (0.71 before): the 50 degree tail and the added drag leave a little less pitch stiffness in the thin air. Borderline, not a loss of control; no law's number was changed to pass it. \\
W1 cruise in turbulence at 20,000 ft, severity 3 (sigma\_w 1.11 m/s): airspeed within 5 kt & Off the design conditions: days, wind, turbulence & As at 60,000 ft: the gust along the heading varies by 15 kt, the airspeed by 11. \\
W2 mission, 25 kt wind, light turbulence: enroute: track within 0.25 nmi for 95 \% of the route & Off the design conditions: days, wind, turbulence & 94 \% against 95 \%, in a 25 kt wind from the runway heading, which is a crosswind or a tailwind on the legs, with light turbulence; phase 7b's own rule (80 \%) passes, and calm air gives 96 \%. The track loop's lookahead is 15 s. \\
\bottomrule
